% Spin-polarized DT neutron source for OpenMC -- figure summary for E. Peterson.
% Compile from spf_prototype/:  pdflatex peterson_figures.tex  (run twice)
\documentclass[11pt]{article}
\usepackage[letterpaper,margin=0.9in]{geometry}
\usepackage{graphicx}
\usepackage{amsmath}
\usepackage{float}
\usepackage[font=small,labelfont=bf]{caption}
\usepackage{microtype}
\usepackage[hidelinks]{hyperref}
\graphicspath{{figs/}}
\setlength{\parskip}{0.4em}
\pagestyle{plain}

\title{\vspace{-1.5em}Spin-Polarized DT Fusion Neutron Source in OpenMC:\\
       Verification and Blanket-Neutronics Results}
\author{}
\date{June 2026}

\begin{document}
\maketitle
\vspace{-2.5em}

\noindent
These figures summarize a self-contained prototype spin-polarized deuterium--tritium (DT)
neutron \emph{source} for OpenMC and its verification. The angular emission follows the
Schwartz (2025) differential cross section $\mathrm{d}\sigma/\mathrm{d}\Omega \propto
\tfrac34 a\,\sin^2\theta + \left(\tfrac23 b + \tfrac13 c\right)\!\left(\tfrac14 +
\tfrac34\cos^2\theta\right)$, with $\theta$ measured from the local magnetic field. The
direction sampler is verified to machine precision against a Python mirror, and the full
OpenMC transport is shown to \emph{recover} an independent analytic neutron wall load
(NWL) computed from Schwartz's closed form and his \texttt{anarr\'ima} package. The model
is then extended to scattering, a tungsten first wall, and a FLiBe breeder blanket. The
work is positioned as \emph{independent validation and a reusable tool}: it cross-validates
the Monte-Carlo spin-polarized-fuel neutronics of Bae et al.\ (\emph{Nucl.\ Fusion}
\textbf{65}, 086051, 2025) through a different source implementation (a compiled
\texttt{openmc::Source}), breeder (FLiBe vs.\ Pb-Li), and geometry. Verification figures
(1--2) use a filamentary/parabolic-ring source in a purely toroidal field with scattering
disabled; quantities are per source neutron unless noted.

\clearpage

\begin{figure}[H]
  \centering
  \includegraphics[width=0.95\textwidth]{tier3_analytic_vs_numerical.png}
  \caption{OpenMC Monte-Carlo transport (markers) reproduces the independent Schwartz
  analytic free-streaming wall load (lines). Each panel is the per-bin ratio of neutron
  wall load for a pure polarization mode to the isotropic case, along one wall of a
  square-cross-section torus (parabolic ring plasma, toroidal field, no scattering). The
  A mode ($\sin^2\theta$) raises the inboard load ${\sim}41\%$ and lowers the outboard
  ${\sim}21\%$; the B/C mode ($1+3\cos^2\theta$) is the mirror image. Bottom-right: the
  standardized residual $(\mathrm{OpenMC}-\mathrm{analytic})/\sigma$ over all walls and
  modes follows $N(0,1)$ ($\mu=+0.02$, $\sigma=0.69$), so the code matches the analytic
  directionality with no detectable bias.}
\end{figure}

\clearpage

\begin{figure}[H]
  \centering
  \includegraphics[width=0.93\textwidth]{tier3_torus3d.png}
  \caption{Computed free-streaming wall load on the four walls of the square-cross-section
  torus, non-polarized (left) versus polarized B/C (right) at constant fusion rate (shared
  colour scale, blue low to red high). This is a render of the tallied load, not a CAD
  geometry. The B/C mode steers neutrons off the inboard center stack toward the outboard
  wall---the benefit relevant to space-constrained spherical-tokamak center stacks.}
\end{figure}

\vspace{1em}

\begin{figure}[H]
  \centering
  \includegraphics[width=\textwidth]{tier5_heating_spectrum.png}
  \caption{Scattering with real materials (cavity, W first wall, steel, Be multiplier,
  FLiBe breeder). Left: energy deposition per source neutron by layer; the FLiBe blanket
  takes ${\sim}91\%$ and the 1\,mm tungsten first wall only ${\sim}0.1\%$. Right: the FLiBe
  flux spectrum shows the 14.1\,MeV source line plus a large down-scattered tail to keV
  energies, the signature of active scattering that drives low-energy $^{6}$Li breeding and
  material damage. Shown for the unpolarized (isotropic) reference: these are bulk-blanket
  quantities, essentially independent of polarization (the surrounding blanket fixes the
  radial attenuation, while polarization only redirects where neutrons land, Figs.~1--2).
  The per-layer heating split varies $<0.3$ percentage points across the A/B/C modes
  (verified), so the unpolarized case is representative.}
\end{figure}

\clearpage

\begin{figure}[H]
  \centering
  \includegraphics[width=\textwidth]{tier6_tbr.png}
  \caption{Tritium breeding in the FLiBe blanket (TBR $=$ tritium per source neutron).
  Left: TBR versus $^{6}$Li enrichment for the unpolarized reference (the optimum is a
  blanket-material property, independent of polarization) rises then plateaus near
  30--50\%, reproducing the published optimum; the 30\% baseline gives TBR $=1.15$
  (ARC/LIBRA band, DEMO goal). Right: TBR for the polarized modes A, B, C versus the
  unpolarized case at constant fusion rate is only modestly polarization-dependent
  ($\pm0.6\%$; B/C higher, A lower), the same sign as Bae et al.\ (2025) ($+2.7\%$ parallel,
  Pb-Li spherical tokamak). Steering to protect the center stack costs little breeding.
  Error bars are Monte-Carlo $1\sigma$.}
\end{figure}

\vspace{1em}

\begin{figure}[H]
  \centering
  \includegraphics[width=0.78\textwidth]{tier7_absolute.png}
  \caption{Absolute (un-normalized) effect with the reaction-rate enhancement restored. At
  fixed fuel density the neutron yield scales with $\eta = a + \tfrac23 b + \tfrac13 c$, so
  A produces $+50\%$ neutrons and C $-50\%$. Bars: fusion rate, absolute tritium production,
  and absolute center-stack load, each relative to unpolarized fuel. A is a power play
  ($+50\%$ power and tritium) but compounds the center-stack load to $1.7\times$; B buys
  steering ($0.85\times$ center stack) at no rate or breeding penalty; C halves everything.
  The TBR per neutron is mode-independent, so self-sufficiency is preserved.}
\end{figure}

\end{document}
